\subsection{\texorpdfstring{\(\mathbf{C}\) 上的拓扑向量空间}{复数域上的拓扑向量空间}}

设 E 是复数域 C 上的一个拓扑向量空间；E 的拓扑也与把标量域限制到 R 所得的 R 上向量空间结构相容。我们用 \(E_{0}\) 表示 E 底下的 R 上拓扑向量空间（I, p.~2）。注意，在 \(E_{0}\) 中，映射 \(x \mapsto ix\)（它不是一个位似）是 \(E_{0}\) 的拓扑向量结构的一个自同构 u，满足 \(u^{2}(x) = -x\)。

反之，设 \(\mathbf{F}\) 是 \(\mathbf{R}\) 上的一个拓扑向量空间，并假设存在一个自同构 \(u\)（\(\mathbf{F}\) 上的），使 \(u^2 = -1_{\mathbf{F}}(1_{\mathbf{F}}\) 是 \(\mathbf{F}\) 的恒等自同构。我们知道（A, IX, § 3, No.~2），这时可以在 \(\mathbf{F}\) 上定义一个关于 \(\mathbf{C}\) 的向量空间结构：写 \(\lambda x = \alpha x + \beta u(x)\)，它对所有 \(\lambda = \alpha + i\beta \in \mathbf{C}\) 与所有 \(x \in \mathbf{F}\) 成立。此外，由于映射 \((\alpha, \beta, x) \mapsto \alpha x + \beta u(x)\)（从 \(\mathbf{R}^2 \times \mathbf{F}\) 到 \(\mathbf{F}\) 中）是连续的，\(\mathbf{F}\) 的拓扑与上面定义的关于 \(\mathbf{C}\) 的向量空间结构相容；若用 \(\mathbf{E}\) 表示这样定义的 \(\mathbf{C}\) 上的拓扑向量空间，则 \(\mathbf{F}\) 就是 \(\mathbf{R}\) 上那个作为 \(\mathbf{E}\) 之底的拓扑向量空间。

注. --- 给定 R 上的一个拓扑向量空间 F，并非总有平方为 \(-1_{F}\) 的自同构 u 存在；例如，在有限奇数维的 R 上向量空间上，不可能定义关于 C 的向量空间结构。

设 E 是 C 上的一个拓扑向量空间，\(E_{0}\) 是 E 底下的 R 上拓扑向量空间。E 中的每个线性流形 M 也是 \(E_{0}\) 中的线性流形，但反之不成立。为避免混淆，我们称关于 C（相应地，关于 R）的向量空间结构的线性流形为复（相应地，实）线性流形。有限维 n（相应地，有限余维 n）的复线性流形是维数 2n（相应地，余维数 2n）的实线性流形。为使 E 的实向量子空间 M 同时也是复向量子空间，必要且充分的是 \(iM \subset M\)。

回忆一下，若 E 与 F 是 C 上的两个拓扑向量空间，则称 E 到 F 中的一个映射为 C-线性的（相应地，R-线性的），如果它对 E 与 F 的关于 C（相应地，R）的向量空间结构是线性映射；每个 C-线性映射显然是 R-线性的，但反之不成立。我们称 E 上的 C-线性型为复线性型，称 E 上的 R-线性型（即 \(E_{0}\) 上的线性型）为实线性型。若 f 是 E 上的一个复线性型，则显然 f 的实部 g = Rf 与虚部 h = If 都是实线性型；此外，关系 \(f(ix) = if(x)\) 蕴含恒等式 \(h(x) = -g(ix)\)；换言之，我们有

\[
f (x) = (\mathscr {R} f) (x) - i (\mathscr {R} f) (i x).\tag{1}
\]

反之，若 g 是 E 上的一个实线性型，则 \(f(x) = g(x) - ig(ix)\) 是 E 上使 Rf = g 的唯一的复线性型；而 f 连续当且仅当 g 连续。

现在设 H 是 E 中的一个复超平面，其方程为 \(f(x) = \alpha + i\beta\)，其中 f 是 E 上的一个复线性型；令 \(g = Rf\)，我们看出 H 是两个实超平面 \(H_{1}\)、\(H_{2}\) 的交，它们的方程分别为 \(g(x) = \alpha\) 与 \(g(ix) = -\beta\)；若 H 是闭的，则 \(H_{1}\) 与 \(H_{2}\) 也是闭的（I, p.~13, 定理 1）。反之，设 \(H_{0}\) 是一个齐次实超平面，其方程为 \(g(x) = 0\)（其中 g 是 E 上的一个实线性型）；则 H——即 \(H_{0}\) 与 \(iH_{0}\) 的交——是一个齐次复超平面，而若 f 是使 Rf = g 的那个复线性型，则 \(f(x) = 0\) 便是 H 的方程；若 \(H_{0}\) 是闭的，则 H 也是闭的。

设 G 是 R 上的一个拓扑向量空间，并设 \(\mathbf{G}_{(\mathbf{C})}\) 是把标量域扩张到 C 后由 G 得到的 C 上向量空间（A, II, § 5.1）。把 G 等同于 \(\mathbf{G}_{(\mathbf{C})}\) 的一个子集，借助映射 \(x \mapsto 1 \otimes x\)。于是，R-线性映射 \((x, y) \mapsto x + i \cdot y\) 是 \(G \times G\) 到 \(\mathbf{G}_{(\mathbf{C})}\) 上的一个双射，借助它，我们把 \(G \times G\) 的乘积拓扑搬到 \(\mathbf{G}_{(\mathbf{C})}\) 上。这样，赋有此拓扑的 \(\mathbf{G}_{(\mathbf{C})}\) 是 C 上的一个拓扑向量空间。我们称 \(\mathbf{G}_{(\mathbf{C})}\) 为 G 的复化的拓扑向量空间。

